% Muc 7.6  --  phu trach: xem report/PHANCONG.md
% nguon: docs/quartus/cnn1d.*.summary (7x7, 21/11/2026); docs/quartus/6x8/, docs/quartus/3x3/ (28/11/2026, gom raw log board);
%        results/hardware_logs/layer0-5_rec101_93600_beat000_20261124-085811.raw.log (7x7);
%        ban ghi 208: results/2026-11-28_a5fb3f12_board_compare/ (3x3), results/2026-11-24_a5fb3f12_board_compare/ (7x7);
%        kiem lai: python scripts/raw_log_check.py <raw log>; du bao: python scripts/tile_cycle_model.py --shape RxC
%        hinh: report/figures/gen_somang.py
\section{So sánh các kích thước mảng PE}\label{sec:somang}
Kích thước mảng là tham số của RTL (\texttt{ARRAY\_ROWS}, \texttt{ARRAY\_COLS} trong \texttt{de10\_lite\_top.v}).
Bộ điều khiển layer tự chia ô theo kích thước mảng và host đọc kích thước qua lệnh \texttt{INFO}, nên cùng một mã RTL và cùng chương trình host chạy được với nhiều cấu hình.
Mục này so ba cấu hình:
\begin{itemize}
  \item $7\times7$: cấu hình chính, dùng ở các mục trước;
  \item $6\times8$: cấu hình được đề xuất ở Mục~\ref{sec:bandwidth} vì $C_{out}=16,32,64$ chia hết cho $8$, và đã được mô hình dự báo trước khi tổng hợp;
  \item $3\times3$: mảng nhỏ nhất chạy được mọi layer, vì FC1 cần ba hàng PE (Mục~\ref{sec:anhxa}).
\end{itemize}
Ba bản dùng cùng mã RTL và chỉ khác hai tham số trên; bản $7\times7$ biên dịch ngày 21/11/2026, hai bản còn lại ngày 28/11/2026, đều bằng Quartus Prime Lite 18.1.

\subsection{Bảng tổng hợp}
Bảng~\ref{tab:somang} gom các số của ba cấu hình.
Mỗi số đi kèm cách thu được: tài nguyên và định thời là báo cáo của Quartus, chu kỳ là số đo trên board, công suất và năng lượng là ước lượng.

\begin{table}[htbp]\centering
\caption[So sánh ba kích thước mảng PE]{So sánh ba kích thước mảng PE trên MAX~10 10M50 ở $50$~MHz.
Chu kỳ đo trên board với nhịp 0 của đoạn \texttt{rec101\_93600}.}\label{tab:somang}
\small\setlength{\tabcolsep}{5pt}
\begin{tabular}{@{}lrrr@{}}
\toprule
Chỉ số & $3\times3$ & $6\times8$ & $7\times7$ \\ \midrule
Số PE & 9 & 48 & 49 \\
\multicolumn{4}{@{}l}{\emph{Báo cáo Quartus}} \\
\quad Phần tử logic (LE) & 18\,683 (38\,\%) & 43\,221 (87\,\%) & 41\,353 (83\,\%) \\
\quad Thanh ghi & 8\,944 & 15\,612 & 15\,697 \\
\quad Bit bộ nhớ & 386\,560 & 451\,871 & 452\,390 \\
\quad Nhân nhúng 9 bit & 8 & 8 & 8 \\
\quad Fmax, mô hình chậm 85\,$^\circ$C (MHz) & 55{,}42 & 54{,}83 & 50{,}21 \\
\quad Độ dư setup nhỏ nhất (ns) & $+1{,}955$ & $+1{,}762$ & $+0{,}082$ \\
\multicolumn{4}{@{}l}{\emph{Đo trên board}} \\
\quad Chu kỳ mỗi nhịp & 1\,097\,565 & 347\,003 & 367\,091 \\
\quad Thời gian tính mỗi nhịp (ms) & 21{,}95 & 6{,}94 & 7{,}34 \\
\quad Hiệu suất PE $\eta$ (\%) & 16{,}29 & 9{,}66 & 8{,}95 \\
\quad Mô hình chu kỳ, lệch so với đo & $-0{,}03\,\%$ & $-0{,}17\,\%$ & $-0{,}16\,\%$ \\
\quad Sáu layer khớp golden & có & có & có \\
\multicolumn{4}{@{}l}{\emph{Ước lượng}} \\
\quad Công suất tổng (mW) & 210{,}19 & 360{,}10 & 332{,}17 \\
\quad Năng lượng tính mỗi nhịp (mJ) & 4{,}61 & 2{,}50 & 2{,}44 \\ \bottomrule
\end{tabular}

\smallskip
\footnotesize $\eta$ tính theo công thức~\eqref{eq:eta} với $R\times C$ của từng cấu hình.
Năng lượng $=$ công suất tổng $\times$ thời gian tính.
Công suất bản $3\times3$ lấy từ Power Analyzer của lần biên dịch ngày 27/11/2026, có số Fitter và định thời trùng lần biên dịch ngày 28/11/2026.
Nguồn: \texttt{docs/quartus/}, \texttt{docs/quartus/6x8/}, \texttt{docs/quartus/3x3/}.
\end{table}

Hình~\ref{fig:somangchiso} vẽ lại sáu chỉ số chính, mỗi chỉ số một trục riêng.

\begin{figure}[htbp]\centering
  \includegraphics[width=\textwidth]{ch07_thuc_nghiem/somang_chiso}
  \caption[Các chỉ số chính của ba kích thước mảng PE]{Các chỉ số chính của ba kích thước mảng PE.
  Đường nét đứt ở ô Fmax là tần số yêu cầu $50$~MHz.}\label{fig:somangchiso}
\end{figure}

\subsection{Thời gian tính theo từng layer}
Hình~\ref{fig:somangchuky} cho số chu kỳ đo được của từng layer.

\begin{figure}[htbp]\centering
  \includegraphics[width=0.92\textwidth]{ch07_thuc_nghiem/somang_chuky}
  \caption{Số chu kỳ từng layer của ba kích thước mảng, đo trên board.}\label{fig:somangchuky}
\end{figure}

\textbf{$6\times8$ so với $7\times7$.}
Bản $6\times8$ cần ít hơn $5{,}5\,\%$ chu kỳ dù ít hơn một PE.
Ở conv0--conv2, các layer có pool dùng $6$ hàng ở cả hai cấu hình, còn $8$ cột chia hết $C_{out}$ nên số ô giảm từ $180$, $150$, $150$ xuống $120$ mỗi layer; số chu kỳ giảm $19{,}6\,\%$, $8{,}0\,\%$ và $9{,}0\,\%$.
conv3 thì chậm hơn $3{,}3\,\%$: layer này không có pool nên bản $7\times7$ dùng đủ $7$ hàng, số ô chỉ giảm từ $70$ xuống $64$, trong khi mỗi ô phải nạp $8$ làn thay vì $7$ (thời gian nạp tỉ lệ với $\max(R_u,C_u)\cdot K$).

\textbf{$3\times3$ so với $7\times7$.}
Bản $3\times3$ có ít hơn $5{,}4$ lần số PE nhưng chỉ chậm hơn $3{,}0$ lần, vì hiệu suất PE cao hơn ($16{,}29\,\%$ so với $8{,}95\,\%$).
Điều này khớp với công thức~\eqref{eq:eta-approx}: mỗi ô chỉ nạp $3$ làn, nên giới hạn $1/(M+1)$ là $1/4$ thay vì $1/8$.
Mức lấp đầy kéo hiệu suất xuống dưới giới hạn đó: ở conv0--conv2 mảng chỉ dùng $2$ trong $3$ hàng (layer có pool cần số hàng chẵn), nên $\eta$ của conv1 và conv2 chỉ khoảng $15\,\%$.

\subsection{Kiểm chứng mô hình chu kỳ và tính đúng}
Số chu kỳ của bản $6\times8$ được mô hình dự báo là $346\,404$ (Bảng~\ref{tab:duBao}) trước khi bản này được tổng hợp; số đo là $347\,003$, lệch $0{,}17\,\%$.
Ở cả ba cấu hình, mô hình thấp hơn số đo của conv0--conv3 đúng $69$, $69$, $133$ và $5$ chu kỳ, tức đúng phần ghi vùng đệm và đóng layer không được mô hình hoá (Mục~\ref{sec:bandwidth}).
Như vậy công thức $\delta=20+R+2C_u+N_{\text{out}}$ đúng từng chu kỳ cho mọi ô của các layer tích chập ở $R=3$, $6$ và $7$, không chỉ ở cấu hình dùng để rút ra nó.
Sai số còn lại nằm ở hai layer FC, nơi mô hình dùng hằng số khớp trên bản $7\times7$ (FC1 lệch $-2{,}5\,\%$ ở $3\times3$ và $-10{,}6\,\%$ ở $6\times8$); hai layer này chỉ chiếm dưới $1\,\%$ số chu kỳ nên tổng vẫn lệch dưới $0{,}2\,\%$.

Về tính đúng, host chạy cả sáu layer của cùng một nhịp trên mỗi cấu hình và so với golden sau mỗi layer.
Dữ liệu đọc về của cả sáu layer khớp từng byte với golden ở cả ba cấu hình; kết quả này kiểm lại được từ nhật ký byte UART bằng \texttt{scripts/raw\_log\_check.py}.
Bản $3\times3$ còn được chạy trên $396$ nhịp đầu của bản ghi 208, và nhãn trùng mô hình int8 phần mềm ở cả $396$ nhịp.
Cả hai bản chưa chạy lại $1\,306$ nhịp golden như bản $7\times7$ (Mục~\ref{sec:bitexact}); riêng bản $6\times8$ mới được kiểm trên một nhịp.

\subsection{Đánh đổi}
\begin{itemize}
  \item \textbf{$6\times8$ nhanh hơn một chút nhưng tốn hơn.} Bản này dùng nhiều hơn $1\,868$ LE ($87\,\%$ so với $83\,\%$) và công suất ước lượng cao hơn $8\,\%$, nên năng lượng mỗi nhịp gần như không đổi ($2{,}50$ so với $2{,}44$~mJ).
  Một phần LE tăng thêm có thể do khối hậu xử lý có một đơn vị cho mỗi cột (\texttt{conv1d\_post\_process.v}), tăng từ $7$ lên $8$; báo cáo Fitter theo khối của bản $6\times8$ không được lưu nên chưa tách được.
  Ước lượng LE thô trong \texttt{scripts/tile\_cycle\_model.py} ($14\,040$ LE cho phần ngoài mảng cộng $557$ LE mỗi PE, lấy từ báo cáo Fitter của bản $7\times7$) cho $40\,776$, thấp hơn thực tế $5{,}7\,\%$; với $3\times3$ ước lượng này cho $19\,053$, cao hơn thực tế $2{,}0\,\%$.
  \item \textbf{Độ dư định thời.} Bản $6\times8$ và $3\times3$ còn dư setup gần $2$~ns, trong khi bản $7\times7$ chỉ còn $+0{,}082$~ns.
  Độ dư phụ thuộc kết quả đặt và đi dây của từng lần biên dịch, nên chênh lệch giữa $6\times8$ và $7\times7$ không được quy cho kích thước mảng.
  \item \textbf{$3\times3$ tiết kiệm tài nguyên nhưng tốn năng lượng hơn.} Bản này chỉ dùng $45\,\%$ số LE và $63\,\%$ công suất ước lượng của bản $7\times7$, nhưng tính lâu gấp $3$ lần nên năng lượng mỗi nhịp gấp $1{,}9$ lần.
  Khi chạy bản ghi 208 kèm cửa sổ hiển thị trực tiếp, với cùng chương trình host và cùng chế độ nạp mô hình một lần như Mục~\ref{sec:hieunang}, bản $3\times3$ mất trung bình $0{,}30$~s mỗi nhịp trên $396$ nhịp, so với $0{,}20$~s của bản $7\times7$ trên $2\,946$ nhịp.
  Cả hai vẫn ngắn hơn khoảng RR trung bình $0{,}61$~s của bản ghi này.
  Chênh lệch $0{,}10$~s lớn hơn nhiều so với $14{,}6$~ms tính toán thêm, nên phần lớn chênh lệch nằm ngoài mảng PE; nguyên nhân chưa được tách và chưa đo lại ở chế độ không hiển thị.
  \item \textbf{Chỗ trống trên chip.} Bản $3\times3$ còn khoảng $31\,000$ LE, đủ để thử các cải tiến bộ nạp ở Mục~\ref{sec:bandwidth}; theo mô hình, $3\times3$ với đọc rộng $\times4$ và nạp chồng tính cần $372\,158$ chu kỳ, gần bằng bản $7\times7$ hiện tại (dự báo, chưa có RTL).
\end{itemize}
